\documentclass{article}
\usepackage[margin=1in]{geometry}
\usepackage{mathtools, amsfonts, amsthm, xcolor, listings}

\title{HW 02}
\author{Sam Ly}

\begin{document}
\maketitle

\section*{Exercise 1 [2]}
Collaborators: Sam Ly
\\
Tools and resources used: I used 
\\
\begin{enumerate}
    \item[Exercise 1]
    {\color{red} Points attempted: 2}.
    {\color{blue} Time spent: 11 minutes}.
    {\color{brown} Difficulty: 2/10}.
    {\color{green!50!black} Self-assessed score: 2}.
    {\color{magenta} Questions/Feedback: None.}

    \item[Exercise 2]
    {\color{red} Points attempted: 1}.
    {\color{blue} Time spent: 11 minutes}.
    {\color{brown} Difficulty: 2/10}.
    {\color{green!50!black} Self-assessed score: 1}.
    {\color{magenta} Questions/Feedback: None.}

    \item[Exercise 3]
    {\color{red} Points attempted: 7}.
    {\color{blue} Time spent: 1 hour + 27 minutes = 87 minutes}.
    {\color{brown} Difficulty: 4/10}.
    {\color{green!50!black} Self-assessed score: 7}.
    {\color{magenta} Questions/Feedback: mathematical correctness.}
\end{enumerate}
\pagebreak
    
\section*{Exercise 2 [1]}
Done! 

I made it to example 1.1.15, which discusses product groups. I was thinking 
about what is required for a product group to be valid. For example, if two 
groups G and H have distinct orders, can they still form a product group? What 
if the order of one group is a multiple/shares factors with the other? What if 
the groups are "vastly different" (in the sense that one group may be finite, 
while another is infinite)? Are there any special properties that emerge from 
products of groups that are not present in the individual groups themselves?
\pagebreak

\section*{Exercise 3 [7]}
\begin{enumerate}
    \item {
        Use the principle of mathematical induction to prove that for all \(k \geq 1\),
        \[(g^k)^{-1} = (g^{-1})^k.\]

        \begin{proof}
            We procede via induction. We first establish base cases by checking 
            that the proposition holds for small \(1 \le k \le 3\). 

            For \(k = 1\), we simply have the expression \((g^1)^{-1} = (g^{-1})^1\),
            which simplifies to \(g^{-1} = g^{-1}\). 

            For \(k = 2\), we have the expression \((g^2)^{-1} = (g^{-1})^2\).
            For clarity, we write \((g \times g)^{-1} = (g^{-1} \times g^{-1})\).

            Now, we can apply the \emph{Inverse of a Product} property which states 
            that \[ (a \times b)^{-1} = b^{-1} \times a^{-1},\]
            to arrive at the conclusion 
            \begin{align*}
                (g^2)^{-1} &= (g^{-1})^2 \\
                (g \times g)^{-1} &= (g^{-1} \times g^{-1}) \\
                g^{-1} \times g^{-1} &= g^{-1} \times g^{-1}.
            \end{align*}

            Proving this for \(k = 3\) requires a little bit more algebraic 
            manipulation. However, the technique used for \(k = 3\) will come in 
            handy for proving the general case. 

            To begin, we write
            \begin{align*}
                (g^3)^{-1} &= (g^{-1})^3 \\
                (g \times g \times g)^{-1} &= g^{-1} \times g^{-1} \times g^{-1}. \\
            \end{align*}
            By \emph{associativity},
            \begin{align*}
                (g \times g \times g)^{-1} &= g^{-1} \times g^{-1} \times g^{-1} \\
                (g \times (g \times g))^{-1} &= g^{-1} \times g^{-1} \times g^{-1}. \\
            \end{align*}

            By reapplying the \emph{Inverse of Product} property, we write
            \begin{align*}
                (g \times g)^{-1} \times g^{-1} &= g^{-1} \times g^{-1} \times g^{-1} \\
                g^{-1} \times g^{-1} \times g^{-1} &= g^{-1} \times g^{-1} \times g^{-1}. \\
            \end{align*}

            Now, we can make our inductive hypothesis: That \((g^k)^{-1} = (g^{-1})^k\)
            for \(k \le n\). 
            
            To prove that our proposition holds for \(k +1\), we use 
            \emph{associativity} to ``pull appart''the LHS expression to get 
            \begin{align*}
                (g^{k+1})^{-1} &= (g^{-1})^{k+1} \\
                (g \times g^k)^{-1} &= (g^{-1})^{k+1} \\
                (g^k)^{-1} \times g^{-1} &= (g^{-1})^{k+1} \\
                (g^{-1})^k \times g^{-1} &= (g^{-1})^{k+1} \\
                (g^{-1})^{k+1} &= (g^{-1})^{k+1}.
            \end{align*}

            Therefore, for all \(k \geq 1\),
            \[(g^k)^{-1} = (g^{-1})^k.\]
        \end{proof}
    }

    \item {
        Use the principle of mathematical induction to prove that for all \(n \geq 1\),
      \[         (g_1 * g_2 * \dots * g_n)^{-1} = g_n^{-1} * \dots * g_2^{-1} * g_1^{-1}.      \]

      \begin{proof}
        We procede via induction. To establish our base cases, we prove our proposition
        for \(1 \le n \le 3\). 

        For \(n = 1\), we have \((g_1)^{-1} = g_1^{-1}\).
        For \(n = 2\), we use the \emph{Inverse of Product} property to see 
        \[ (g_1 * g_2)^{-1} = g_2^{-1} * g_1^{-1}. \]

        For \(n = 3\), we use similar reasoning as the previous problem.
        \begin{align*}
            (g_1 * g_2 * g_3)^{-1} = g_3^{-1} * g_2^{-1} * g_1^{-1} \\
            (g_1 (* g_2 * g_3))^{-1} = g_3^{-1} * g_2^{-1} * g_1^{-1} \\
            (g_2 * g_3)^{-1} * (g_1 )^{-1} = g_3^{-1} * g_2^{-1} * g_1^{-1} \\
            g_3^{-1} * g_2^{-1} * g_1^{-1} = g_3^{-1} * g_2^{-1} * g_1^{-1}. \\
        \end{align*}

        As an inductive hypothesis, we say 
        \[         (g_1 * g_2 * \dots * g_n)^{-1} = g_n^{-1} * \dots * g_2^{-1} * g_1^{-1}.      \]
        for \(n \le k\). 

        Thus, we have 
        \[         (g_1 * g_2 * \dots * g_k)^{-1} = g_k^{-1} * \dots * g_2^{-1} * g_1^{-1}.      \]

        We can show that our proposition holds for \(n = k+1\) with
        \begin{align*}
            (g_1 * g_2 * \dots * g_k * g_{k+1})^{-1}\\
            ((g_1 * g_2 * \dots * g_k) * g_{k+1})^{-1}\\
            (g_{k+1})^{-1} * (g_1 * g_2 * \dots * g_k)^{-1}\\
            g_{k+1}^{-1} * g_k^{-1} * \dots * g_2^{-1} * g_1^{-1}
        \end{align*}

        Therefore, for all \(n \geq 1\),
      \[         (g_1 * g_2 * \dots * g_n)^{-1} = g_n^{-1} * \dots * g_2^{-1} * g_1^{-1}.      \]

      \end{proof}
    }

    \item {
        Give an example of a group \((G,\star)\) along with two elements \(g_1, g_2 \in G\) such that \((g_1 \star g_2)^{-1} \neq g_1^{-1} \star g_2^{-1}\).
    
        This statement applies to all groups \(G\) such that the elements of 
        \(g \in G\) do not commute. In other words, all \emph{non-abelian} groups.
        One example of a \emph{non-abelian} group are the dihedral group \(D_n\). 
    }

    \item 
        Prove that function composition is associative.

        \begin{proof}
            Suppose we have functions \(f, g, h\). We say that the composition
            \(f \circ g \circ h\) is ``\(f\) after \(g\) after \(h\)''.

            However, we can see that \(f \circ (g \circ h)\) is just ``
            \(f\) after (\(g\) after \(h\))''. Notice that the groupings don't 
            actually change the order of function applications!

            Similarly, \((f \circ g) \circ h\) also doesn't change the order 
            of function applications. 

            Therefore, function composition must be associative.
        \end{proof}
    
    \item 
        Prove that the set of bijections \(\mathbb{R} \to \mathbb{R}\) (i.e. 
        those whose graphs pass the "horizontal line test") form a group under 
        function composition \(\circ\).
    
        \begin{proof}
            To show that the set of bijections \(\mathbb{R} \to \mathbb{R}\) form 
            a group under function composition \(\circ\), we must show that such 
            as set exhibit the group axioms. 

            First, a group must have an identity element \(e\) such that 
            \(e \circ f = f \circ e = f\). We see that \(e(x) = x\) fits the bill.

            Then, all elements of a group must have an inverse. By the definition 
            of bijections, our functions must be invertable. Thus, for all 
            functions \(f \in G\), there exists a function \(f^{-1} \in G\) such 
            that \(f \circ f^{-1} = f^{-1} \circ f = e\). 

            Our binary operation is the function composition. We see that by 
            definition of bijections, the composition of two bijection is another 
            bijection. Thus, our functions are closed under composition. 

            Lastly, we must show that our set of bijections is associative under 
            function composition. We showed this in the previous problem.

            Therefore, set of bijections \(\mathbb{R} \to \mathbb{R}\) form a 
            group under function composition.

        \end{proof}

    \item
        Prove or disprove that \(\{7n \mid n \in \mathbb{Z}\}\) forms a group under addition. We sometimes call this set \(7\mathbb{Z}\).
    
        \begin{proof}
            Let \(S = \{7n \mid n \in \mathbb{Z}\} \).
            First, we see that 0 is our identity element under addition, and 
            \(0 \in S\).

            Then, we see that all elements have an inverse that is also in \(S\). 
            For any element \(a = 7n\), \(a^{-1} = -7n\).

            Our binary operation is also closed because for any \(a = 7n, b = 7m\),
            \(a + b = 7n + 7m = 7(n + m) \in S\).

            Lastly, we see that our operation is associative by definition of addition.

            Therefore, \(S\) forms a group under addition.
        \end{proof}

    \item  
        Prove or disprove the set of matrices of the form
      \[         \left\{           \begin{pmatrix}1&x\\0&1\end{pmatrix}           \middle| x \in \mathbb{R}         \right\}       \] forms a group under matrix multiplication.
    
        \begin{proof}
            First, we see that \(I\) (identity matrix) is in our set, so we have a valid identity
            element. 

            Then, we see that for all \(x \in \mathbb{R}\), the matrix \(\begin{pmatrix}1&x\\0&1\end{pmatrix}\)is invertable.
            Specifically,
            \[\begin{pmatrix}1&x\\0&1\end{pmatrix}^{-1} = \begin{pmatrix}1&-x\\0&1\end{pmatrix}.\]

            Our binary operation of matrix multiplication is also closed because
            for \(x, y \in \mathbb{R}\). 
            \[
                \begin{pmatrix}1&x\\0&1\end{pmatrix} \begin{pmatrix}1&y\\0&1\end{pmatrix} = \begin{pmatrix}1&x+y\\0&1\end{pmatrix}.
            \]
            Since \(x + y \in \mathbb{R}\), the new matrix exists within our set.

            Lastly, by definition of matrix multiplication, our operation must 
            also be associative.

            Therefore, set of matrices of the form
      \[         \left\{           \begin{pmatrix}1&x\\0&1\end{pmatrix}           \middle| x \in \mathbb{R}         \right\}       \] forms a group under matrix multiplication.
            form a group under matrix multiplication.

        \end{proof}
    
\end{enumerate}
\pagebreak

\section*{title}

\end{document}
